Mais importante que saber programar é saber encontrar falhas no próprio código. Por isso, o processo 
de depuração (debugging) é tão crucial. Não serão poucas as vezes em que, durante a criação de um software 
mais complexo, você se deparará com erros que parecem não fazer sentido. Seja um novato ou alguém com décadas 
de experiência, todos enfrentarão esse tipo de problema.

Depurar é uma arte à parte, algo que está em constante evolução dentro do universo da programação. Existem 
inúmeras ferramentas e estratégias que podem ser utilizadas nesse processo; porém, serão abordadas aqui duas 
das ferramentas mais conhecidas, disponíveis gratuitamente e com amplo material de suporte na internet:

\begin{itemize}
    \item \textbf{GDB}: uma ferramenta TUI (Text-based User Interface) que permite depurar aplicações em C 
        de maneira detalhada por meio de comandos simples;
    \item \textbf{Valgrind}: outra ferramenta TUI, usada principalmente para detectar vazamentos de memória 
        em programas escritos em C.
\end{itemize}

Antes de ensinar como utilizar essas duas ferramentas, vamos precisar definir um código exemplo, que deve 
conter alguns erros a serem consertados. Portanto, vamos considerar o seguinte código:

\begin{excode}
    #include <stdio.h>
    #include <stdlib.h>

    int main()
    {
        int* x = malloc(sizeof(*x));
        *x = 0;

        while (*x < 10) {
            printf("X é menor que 10!\n");
            x += 1;
        }

        return 0;
    }
\end{excode}

O objetivo do código acima é simples: alocar memória para um inteiro \texttt{x}, armazenar o valor 0 
nessa posição, incrementar esse valor de 1 em 1 enquanto ele for menor que 10 e exibir um texto na tela a cada 
iteração. Nesse ponto, a execução da aplicação deve ser finalizada.

Quando compilado, o código não gera \textbf{erros de compilação}, apenas alguns \textit{warnings} 
(avisos) que podem ser ignorados dependendo das configurações. Porém, ao executá-lo, você vai ter contato 
ao que chamamos de \textbf{undefined behavior} (ou comportamento indefinido) pois não é possível dizer como o 
código vai se comportar, podendo acontecer qualquer uma das seguintes situações:

\begin{itemize}
    \item Executar o \texttt{printf} uma única vez.
    \item Executar o \texttt{printf} várias vezes.
    \item Entrar em um looping infinito.
    \item Código crashar com um erro de segmentação de memória.
\end{itemize}

No meu caso, ao executar o código na minha máquina foi possível obter um erro de segmentação de memória, como 
é mostrado abaixo:

\begin{excode}
    X é menor que 10!
    X é menor que 10!
    X é menor que 10!
    ...
    X é menor que 10!
    Segmentation fault
\end{excode}

Para resolver isso, vamos aprender a utilizar as ferramentas de debugging.

\begin{postscript}
    Vale ressaltar que para utilizar ferramentas de debugging de maneira mais prática e 
    eficiente, é necessário que seja utilizada a flag \texttt{-g} durante o processo de 
    compilação de uma aplicação. Por exemplo:

    \begin{textbox}
        gcc -o main main.c -g 
    \end{textbox}

    Usar essa opção faz com que o programa seja compilado com as informações extras de 
    depuração que facilitam MUITO o uso de ferramentas como o GDB ou Valgrind. Nos exemplos 
    abaixo vamos levar em consideração que você compilou o programa corretamente usando a flag \texttt{-g}
\end{postscript}

\subsection{GDB}

A princípio o GDB acaba sendo fundamental por ser uma ferramenta que permite controlar a execução 
de um programa, possibilitando inspecionar o código linha a linha, definir breakpoints (pontos de 
interrupção), verificar o estado das variáveis, entre outras funcionalidades.

Para entender seu funcionamento, vamos aprender alguns comandos básicos. Para iniciar a depuração, 
é necessário abrir o programa pelo GDB com o seguinte comando:

\begin{textbox}
    gdb [nome_do_programa]
\end{textbox}

Após executar esse comando, o GDB carrega o programa em um contexto de depuração, sem executá-lo 
imediatamente. Se tudo ocorrer corretamente, você deve ver no terminal:

\begin{textbox}
    Reading symbols from main...
    (gdb)
\end{textbox}

A partir daí, é possível definir pontos de interrupção para investigar possíveis erros. Para visualizar 
o código-fonte, utiliza-se o comando:

\begin{textbox}
    list [nome_da_funcao]
    // ou
    l [nome_da_funcao]
\end{textbox}

Este comando exibe o código da função especificada com suas respectivas numerações de linha. Em 
nosso caso, usaremos para visualizar a função \texttt{main}.

\begin{postscript}
    Vale mencionar que o comando \texttt{list} exibe uma quantidade limitada de linhas por vez. 
    Para ver as próximas linhas, basta digitar o comando novamente ou pressionar Enter. Ao chegar 
    ao final da função, uma mensagem de aviso será exibida.
\end{postscript}

\begin{textbox}
    (gdb) l main
    1       #include <stdio.h>
    2       #include <stdlib.h>
    3
    4       int main()
    5       {
    6           int* x = malloc(sizeof(*x));
    7           *x = 0;
    8
    9           while (*x < 10) {
    10              printf("X é menor que 10!\n");
    (gdb) 
    11              x += 1;
    12          }
    13
    14          return 0;
    15      }
    (gdb) 
    Line number 16 out of range; main.c has 15 lines.
\end{textbox}

Tendo o número de cada linha de código é possível definir os pontos de interrupção, neste caso, 
vamos definir um ponto na função \texttt{main} com o seguinte comando:

\begin{exterminal}
    break main
    // ou
    b main
\end{exterminal}

Este comando cria um \textbf{breakpoint} no início da função especificada e exibe uma saída semelhante a:

\begin{exterminal}
    (gdb) break main
    Breakpoint 1 at 0x1151: file main.c, line 6.
\end{exterminal}

\begin{postscript}
    No programa de exemplo não há muitas funções para depurar, porém, em códigos maiores e/ou mais complexos, 
    bons pontos de interrupção podem ser definidos ao responder uma pergunta simples: "Onde o código começou 
    a apresentar problemas?".

    No fim das contas, trata-se de um trabalho investigativo (e por vezes cansativo), mas com a prática você 
    desenvolverá uma noção de onde e como procurar.
\end{postscript}

Em seguida, inicia-se o processo de depuração com o comando \texttt{run} (ou \texttt{r}):

\begin{exterminal}
    (gdb) r
    Starting program: /file/path/main 
    [Thread debugging using libthread_db enabled]
    Using host libthread_db library "/lib/x86_64-linux-gnu/libthread_db.so.1".

    Breakpoint 1, main () at main.c:6
    6           int* x = malloc(sizeof(*x));
\end{exterminal}

Nessa saída, é possível observar informações como:
\begin{itemize}
    \item O caminho do programa que está sendo depurado
    \item O identificador do \textbf{breakpoint} ativo
    \item A função em execução
    \item A próxima linha de código a ser executada
\end{itemize}

Neste momento, a linha 6 ainda não foi executada, portanto não há nada a verificar. Para executar a linha 
atual e avançar para a próxima, utiliza-se o comando \texttt{next} (ou \texttt{n}):

\begin{exterminal}
    (gdb) n
    7           *x = 0;
    (gdb) n
    9           while (*x < 10)
\end{exterminal}

Para verificar o valor de variáveis vamos usar o seguinte comando:

\begin{exterminal}
    print nome_da_variavel
    // ou
    p nome_da_variavel
\end{exterminal}

Como executamos as linhas que alocam memória para a variável \texttt{x} usando a função \texttt{malloc} 
e definimos o seu valor para 0, podemos verificar se as operações foram realizadas corretamente:

\begin{exterminal}
    (gdb) p x
    $3 = (int *) 0x5555555592a0
    (gdb) p *x
    $4 = 0
\end{exterminal}

Note que foi necessário desreferenciar o ponteiro para acessar o valor armazenado na memória alocada. Sem a 
desreferência, o GDB exibe apenas o endereço da variável em formato hexadecimal.

Até este ponto, nada de anormal foi detectado. Portanto, vamos continuar com a 
depuração passo a passo, usando apenas os comandos abordados acima:

\begin{exterminal}
    (gdb) n
    9           while (*x < 10) {
    (gdb) p x
    $1 = (int *) 0x5555555592a0
    (gdb) p *x
    $2 = 0
    (gdb) n
    10              printf("X é menor que 10!\n");
    (gdb) n
    X é menor que 10! // Programa exibindo o output de printf
    11              x += 1;
    (gdb) n
    9           while (*x < 10) {
    (gdb) p x
    $3 = (int *) 0x5555555592a4 // Ponteiro foi incrementado em vez do valor!!!!
    (gdb) p *x
    $4 = 0
\end{exterminal}

Como é possível observar no comentário, uma simples verificação das operações revelou que o comportamento 
inesperado do código era causado por um pequeno erro na linha 11:

\begin{excode}
    x += 1;  // Incrementando o ponteiro
    // Deve ser corrigido para
    *x += 1; // Incrementando o valor
\end{excode}

Após identificar o erro, podemos encerrar a depuração com o comando \texttt{quit}:

\begin{exterminal}
    (gdb) quit
    A debugging session is active.

            Inferior 1 [process 19336] will be killed.

    Quit anyway? (y or n) 
\end{exterminal}

Com a correção aplicada e o comportamento indefinido eliminado, ao executar o programa, obtemos o resultado esperado:

\begin{exterminal}
    \$ ./main
    X é menor que 10!
    X é menor que 10!
    X é menor que 10!
    X é menor que 10!
    X é menor que 10!
    X é menor que 10!
    X é menor que 10!
    X é menor que 10!
    X é menor que 10!
    X é menor que 10!
\end{exterminal}

Com as informações desta seção, é possível perceber que, com apenas alguns comandos básicos, já é possível realizar 
processos de depuração de maneira relativamente eficaz.

O GDB é uma ferramenta que oferece comandos mais avançados, mas como não é o foco deste livro se aprofundar nela, 
recomendo a você que pesquise mais sobre a ferramenta posteriormente para explorar todo o seu potencial.

\begin{postscript}
    Vale mencionar que, durante os exemplos de depuração acima, repeti várias vezes comandos como \texttt{next}. Entretanto, 
    você pode simplesmente pressionar \textbf{Enter} para que o GDB repita o último comando, facilitando a passagem por 
    trechos mais longos de código.

    Além disso, existem \textit{plugins} que podem ser utilizados em conjunto com o GDB. Um que utilizo bastante é o 
    \textbf{GEF (GDB Enhanced Features)}, que adiciona funcionalidades extras e uma interface melhorada à ferramenta original.
\end{postscript}

\subsection{Valgrind}

O Valgrind, apesar de possuir um funcionamento bem diferente do GDB, é uma ferramenta igualmente fundamental. 
Ele serve como um meio simples e relativamente fácil de usar para detectar vazamentos de memória. O programa funciona 
da seguinte forma: o usuário deve passar algumas flags juntamente com o executável que será verificado; o Valgrind 
então executa o programa e gera um relatório final com informações sobre possíveis problemas encontrados.

Para exemplificar o uso do Valgrind, continuaremos utilizando o código corrigido na seção anterior, como forma de 
complementar as correções de erros que o GDB não nos alerta.

\begin{excode}
    #include <stdio.h>
    #include <stdlib.h>

    int main()
    {
        int* x = malloc(sizeof(*x));
        *x = 0;

        while (*x < 10) {
            printf("X é menor que 10!\n");
            *x += 1;
        }

        return 0;
    }
\end{excode}

À primeira vista, o código aparenta funcionar corretamente, sem qualquer tipo de erro. Isso se deve ao fato de que o 
problema presente no código não é um erro de lógica, mas sim um vazamento de memória. Para não adiantar a resposta, 
vamos tentar descobrir o erro usando o Valgrind.

Para verificar o código por meio do Valgrind, precisamos novamente compilar o programa utilizando a flag \texttt{-g}, 
assim como foi feito na seção anterior para depuração com o GDB:

\begin{exterminal}
    gcc -o main main.c -g
\end{exterminal}

Em seguida, para iniciar a busca por vazamentos de memória, utilizamos um único comando no terminal com algumas flags 
específicas. O Valgrind executará o programa e gerará um relatório ao final:

\begin{exterminal}
    valgrind --tool=memcheck --leak-check=yes --show-reachable=yes --num-callers=20 --track-fds=yes ./main
\end{exterminal}

\subsubsection{Explicação das flags}

\begin{itemize}
    \item \textbf{--tool=memcheck}: Especifica qual ferramenta do Valgrind será utilizada. O \texttt{memcheck} é a 
        ferramenta padrão para verificação de memória e será usada neste exemplo.
    \item \textbf{--leak-check=yes}: Ativa ou desativa a busca por vazamentos de memória ao final da execução. Quando 
        definido como \texttt{yes} ou \texttt{full}, todos os vazamentos individuais são exibidos em detalhes, incluindo 
        sua classificação como possíveis erros.
    \item \textbf{--show-reachable=yes}: Define se devem ser exibidas no relatório as regiões de memória que não foram 
        liberadas ao sistema, mas que ainda possuíam ponteiros apontando para elas ao final da execução do programa.
    \item \textbf{--num-callers=20}: Define a profundidade do \textit{stacktrace} (rastreamento da pilha de chamadas), 
        ou seja, quantas entradas serão exibidas para rastrear os pontos onde os erros ocorreram no código. O valor pode 
        variar de 0 a 500. É importante notar que um número maior de entradas pode tornar a execução mais lenta, pois o 
        programa precisa gerenciar e armazenar mais informações.
    \item \textbf{--track-fds=yes}: Faz com que sejam exibidos, além dos vazamentos de memória, os descritores de arquivo 
        (\textit{file descriptors}) que permaneceram abertos ao final da execução. O relatório inclui informações sobre 
        onde no código o arquivo foi aberto, além de detalhes sobre \textit{sockets} e nomes de arquivo quando disponíveis.
    \item \textbf{./main}: O executável a ser analisado pela ferramenta.
\end{itemize}

Depois de executar o comando acima, o seguinte relatório é exibido no terminal:

\begin{exterminal}
    ==4405== Memcheck, a memory error detector
    ==4405== Copyright (C) 2002-2022, and GNU GPL'd, by Julian Seward et al.
    ==4405== Using Valgrind-3.19.0 and LibVEX; rerun with -h for copyright info
    ==4405== Command: ./main
    ==4405== 
    X é menor que 10!
    X é menor que 10!
    X é menor que 10!
    X é menor que 10!
    X é menor que 10!
    X é menor que 10!
    X é menor que 10!
    X é menor que 10!
    X é menor que 10!
    X é menor que 10!
    ==4405== 
    ==4405== FILE DESCRIPTORS: 3 open (3 std) at exit.
    ==4405== 
    ==4405== HEAP SUMMARY:
    ==4405==     in use at exit: 4 bytes in 1 blocks
    ==4405==   total heap usage: 2 allocs, 1 frees, 1,028 bytes allocated
    ==4405== 
    ==4405== 4 bytes in 1 blocks are definitely lost in loss record 1 of 1
    ==4405==    at 0x48417B4: malloc (vg_replace_malloc.c:381)
    ==4405==    by 0x10915A: main (main.c:6)
    ==4405== 
    ==4405== LEAK SUMMARY:
    ==4405==    definitely lost: 4 bytes in 1 blocks
    ==4405==    indirectly lost: 0 bytes in 0 blocks
    ==4405==      possibly lost: 0 bytes in 0 blocks
    ==4405==    still reachable: 0 bytes in 0 blocks
    ==4405==         suppressed: 0 bytes in 0 blocks
    ==4405== 
    ==4405== For lists of detected and suppressed errors, rerun with: -s
    ==4405== ERROR SUMMARY: 1 errors from 1 contexts (suppressed: 0 from 0)
\end{exterminal}

Pelas informações do relatório podemos ver facilmente que, foram perdidos 4 bytes em 1 bloco de memória, que foi alocado na função 
\texttt{main}, especficicamente na linha 6. Portanto, tudo que devemos fazer é verificar na linha 6 em que variável foi guardada a 
o bloco de memória perdido e verificar porquê ele não foi liberado corretamente.

\begin{excode}
    ...
    int* x = malloc(sizeof(*x));
    ...
\end{excode}

No código é possível notar que após alocarmos memória para a variável \texttt{x}, em nenhum outro momento o espaço de memória da 
variável foi liberado de volta ao sistema através da função \texttt{free}, o que ocasiona num problema de vazamento de memória que 
não é informado ao usuário em momento algum durante a execução do programa, mas que em sistemas críticos ou em produtos reais que 
possuem um hardware mais limite, pode ocasionar em problemas graves.

Por fim, tudo que precisamos fazer para solucionar o caso acima, é usar devidamente a função \texttt{free} para liberar o espaço 
de memória alocado para \texttt{x} quando ele não for mais ser utilizado pelo programa.

\begin{excode}
    #include <stdio.h>
    #include <stdlib.h>

    int main()
    {
        int* x = malloc(sizeof(*x));
        *x = 0;

        while (*x < 10) {
            printf("X é menor que 10!\n");
            *x += 1;
        }

        free(x); // Libera o espaço de memória alocado para 'x'
        return 0;
    }
\end{excode}

Após fazer a correção exibida no código acima, é possível verificar o código novamente pelo Valgrind utilizando o mesmo comando, e 
teremos o seguinte relatório:

\begin{exterminal}
    ==4669== Memcheck, a memory error detector
    ==4669== Copyright (C) 2002-2022, and GNU GPL\'d, by Julian Seward et al.
    ==4669== Using Valgrind-3.19.0 and LibVEX; rerun with -h for copyright info
    ==4669== Command: ./main
    ==4669== 
    X é menor que 10!
    X é menor que 10!
    X é menor que 10!
    X é menor que 10!
    X é menor que 10!
    X é menor que 10!
    X é menor que 10!
    X é menor que 10!
    X é menor que 10!
    X é menor que 10!
    ==4669== 
    ==4669== FILE DESCRIPTORS: 3 open (3 std) at exit.
    ==4669== 
    ==4669== HEAP SUMMARY:
    ==4669==     in use at exit: 0 bytes in 0 blocks
    ==4669==   total heap usage: 2 allocs, 2 frees, 1,028 bytes allocated
    ==4669== 
    ==4669== All heap blocks were freed -- no leaks are possible
    ==4669== 
    ==4669== For lists of detected and suppressed errors, rerun with: -s
    ==4669== ERROR SUMMARY: 0 errors from 0 contexts (suppressed: 0 from 0)
\end{exterminal}

Como é possível notar lendo o relatório acima, o programa continua executando da mesma forma. Porém, desta vez, sem nenhum 
vazamento de memória, já que todos os espaços de memória alocados foram liberados corretamente.

\subsection{Outras estratégias}

Nesta subseção será abordada uma técnica mais informal de depuração em C, porém ainda bastante utilizada em casos mais simples onde um processo de depuração mais rudimentar seja satisfatório.

Um método que, por mais simples que pareça, pode ser utilizado em alguns casos é o uso de \texttt{printf} para verificar valores de um módulo específico do código, ajudando a identificar se algo está sendo calculado corretamente ou até mesmo para confirmar se determinadas seções estão sendo executadas na ordem esperada.

Considere o seguinte código que soma apenas os números pares de um vetor, mas apresenta um comportamento inesperado:

\begin{excode}
    #include <stdio.h>

    int soma_pares(int* arr, int n) {
        int soma = 0;
        for (int i = 0; i <= n; i++) {
            if (arr[i] % 2 == 0) {
                soma += arr[i];
            }
        }
        return soma;
    }

    int main() {
        int numeros[] = {1, 2, 3, 4, 5, 6};
        int tamanho = 6;
        int resultado = soma_pares(numeros, tamanho);
        printf("Soma dos pares: %d\n", resultado);
        return 0;
    }
\end{excode}

O objetivo do código é simples: percorrer o vetor \texttt{numeros}, acumular apenas os valores pares e exibir o 
resultado. O valor esperado seria \texttt{12} (2 + 4 + 6), porém ao executar o programa, o resultado obtido é diferente:

\begin{exterminal}
    Soma dos pares: 1027435228
\end{exterminal}

O resultado pode variar entre execuções. Esse é um sinal claro de \textbf{undefined behavior}. Caso você tenha suspeitas 
diretas do porquê disso, é possível utilizar o \texttt{printf}, antes de usar ferramentas mais avançadas como o GDB, 
para inspecionar o que está acontecendo a cada iteração do laço:

\begin{excode}
    #include <stdio.h>

    int soma_pares(int* arr, int n) {
        int soma = 0;
        for (int i = 0; i <= n; i++) {
            printf("[debug] i=%d | arr[i]=%d | soma=%d\n", i, arr[i], soma);
            if (arr[i] % 2 == 0) {
                soma += arr[i];
            }
        }
        return soma;
    }

    int main() {
        int numeros[] = {1, 2, 3, 4, 5, 6};
        int tamanho = 6;
        int resultado = soma_pares(numeros, tamanho);
        printf("Soma dos pares: %d\n", resultado);
        return 0;
    }
\end{excode}

Ao executar o código modificado, obtemos uma saída semelhante a:

\begin{exterminal}
    [debug] i=0 | arr[i]=1 | soma=0
    [debug] i=1 | arr[i]=2 | soma=0
    [debug] i=2 | arr[i]=3 | soma=2
    [debug] i=3 | arr[i]=4 | soma=2
    [debug] i=4 | arr[i]=5 | soma=6
    [debug] i=5 | arr[i]=6 | soma=6
    [debug] i=6 | arr[i]=812595920 | soma=12
    Soma dos pares: 812595932
\end{exterminal}

A saída torna o problema evidente, na iteração \texttt{i=6}, o programa acessa \texttt{arr[6]}, que está 
\textbf{fora dos limites do vetor} (cujos índices válidos vão de \texttt{0} a \texttt{5}). O valor lido é lixo de 
memória, e dependendo do ambiente, pode ou não alterar o resultado final, o que explica o comportamento não determinístico.

A causa do problema está na condição do laço:

\begin{excode}
    for (int i = 0; i <= n; i++)
    // Deve ser corrigido para
    for (int i = 0; i < n; i++)
\end{excode}

Após aplicar a correção e remover os \texttt{printf} de depuração, o programa passa a produzir o resultado correto de forma consistente:

\begin{exterminal}
    Soma dos pares: 12
\end{exterminal}

\begin{postscript}
    Vale destacar que os \texttt{printf} inseridos para depuração devem ser removidos do código ao final do processo. Deixá-los no código de produção polui a saída do programa e pode impactar negativamente a performance em laços muito longos.

    Uma prática comum para evitar esse problema é utilizar macros condicionais que ativam as mensagens de depuração apenas quando necessário, como por exemplo:

    \begin{excode}
        #ifdef DEBUG
            printf("[debug] i=%d | arr[i]=%d\n", i, arr[i]);
        #endif
    \end{excode}

    Dessa forma, as mensagens só aparecem quando o programa é compilado com a flag \texttt{-DDEBUG}, mantendo o código de produção limpo.
\end{postscript}
